\section{Why bottlenecks are worth studying}
\label{sec:graphs}

Represent a finite abstraction as a directed graph from initial state $s$ to forbidden state $t$. A set of removable internal vertices that intersects every $s$--$t$ path is a vertex cut. Removing it prevents the forbidden state in that graph. For a finite graph with nonadjacent $s,t$, unit-cost removable internal vertices, and no extra path restrictions, Menger's theorem equates minimum cut size with the maximum number of internally vertex-disjoint paths \cite{menger}. We use this standard result as a design lens, not as a new theorem.

A direct, harmless example makes its value precise. A layered graph with eight fully connected stages of widths $(4,4,4,4,4,4,4,4)$ has $4^8=65{,}536$ paths and minimum vertex cut four. Changing the widths to $(2,8,4,4,4,4,4,4)$ leaves the path count at $65{,}536$ and reduces the minimum cut to two. Cutting that narrow stage blocks every path in the specified graph.

\begin{figure}[t]
\centering
\begin{tikzpicture}[x=.7cm,y=.7cm,font=\footnotesize]
\node[anchor=west,font=\small\bfseries] at (0,3.5) {Same route count; different cut cost};
\foreach \i in {0,...,7}{
  \draw[fill=teal!60,draw=teal] (\i,1.9) rectangle +(0.65,.8);
  \node at (\i+.325,2.3) {4};
}
\node[anchor=west] at (0,1.55) {$65{,}536$ routes; minimum cut $=4$};
\foreach \i/\v in {0/2,1/8,2/4,3/4,4/4,5/4,6/4,7/4}{
  \draw[fill=violet!30,draw=violet] (\i,-.15) rectangle +(0.65,\v*.17);
  \node[anchor=south] at (\i+.325,-.12+\v*.17) {\v};
}
\draw[gold,line width=1.2pt] (-.06,-.22) rectangle (.71,.29);
\node[anchor=west] at (0,-.65) {$65{,}536$ routes; minimum cut $=2$};
\node[anchor=west,font=\scriptsize] at (0,-1.1) {Bars show stage widths; every adjacent pair is fully linked.};
\end{tikzpicture}
\caption{An exact combinatorial example, recomputed by the supplied Rust path counter and vertex-cut solver. No real-world pathway data enter this figure. The narrowest stage matters because all paths traverse it, not because the paths are easy to enumerate.}
\label{fig:cuts}
\end{figure}

Three qualifications prevent overreach. First, a partial graph can have a cheap cut while the real system has an unobserved bypass. Second, real controls have unequal costs, residual leakage, and restrictions on what can be blocked; some plans require conjunctions better represented by hypergraphs. Third, route count can matter to stochastic opportunity and search models even when it does not determine a structural cut. We therefore retain the insight ``look for shared controllable prerequisites'' and reject the universal slogan that pathway count is irrelevant.

This research can begin with permissions, credentials, transaction boundaries, and shared service dependencies in harmless environments. It does not require generating a catalog of catastrophic methods. A cut specification may itself reveal sensitive information, so even abstraction is not a guarantee of harmless disclosure. The near-term program below uses synthetic violations and existing authorized test environments.
